\subsection{Implementation Details}

\paragraph{Supervision and model selection.} JobBERT conditions start from the same released encoder and CRF. Earlier supervision (B1 in the experiment files) uses the historical extraction labels; revised supervision (B2) combines updated guidance, LLM annotation, and adjudication. The conditions share training and development texts, with label changes in 1,440 training records and 107 development records. Development typed exact F1 selects checkpoints, so this comparison changes training supervision and selection targets together.

Qwen starts from its pretrained Instruct weights and fits a fresh LoRA adapter for each of seeds 42, 43, and 44 on the revised 2,150/169 split. Training uses a fixed system instruction, user template, and five schematic examples. Targets contain the literal span text, occurrence index, and type. Loss applies to the assistant JSON and closing turn token. Training uses BF16, and development exact F1 selects the adapter. Table~\ref{tab:training-settings} gives the remaining settings. Historical JobBERT implementation gaps are documented in \hyperref[app:d]{Appendix D}.

\begin{table}[htbp]
\centering
\caption{Original JobBERT and Qwen training settings. Both select checkpoints by development typed exact F1; reproduction notes retain implementation and selection records.}
\label{tab:training-settings}
\small
\begin{tabularx}{\linewidth}{@{}P{0.27\linewidth}P{0.27\linewidth}L@{}}
\toprule
Setting & JobBERT-zh B1/B2 & Qwen LoRA \\
\midrule
Learning rate & $2\times10^{-5}$ & $10^{-4}$ \\
Batch size & 16 & 1; 16-step gradient accumulation \\
Maximum input length & 256 tokens & 8,192 tokens \\
Training limit & 6 epochs; patience 2 & 6 epochs; patience 2 \\
\bottomrule
\end{tabularx}
\par\smallskip\RaggedRight\footnotesize Qwen uses LoRA with rank 16, scaling 32, and dropout 0.05, applied to the attention q/k/v/o projections.
\end{table}

\paragraph{Inference.} Each request contains one sentence's identifier and text, together with the shared instruction and five fixed examples. A parser converts occurrence-based output into character spans. Qwen baseline and adapter inference use the same base checkpoint, prompts, parser, BF16 precision, deterministic decoding, and 2,048-token output cap. DeepSeek uses caps of 2,048 tokens without thinking and 8,192 with thinking. \hyperref[app:c]{\mbox{Appendix~C}} lists service settings and output checks.
